% !TEX root = ../Projektdokumentation.tex
\section{Fazit}
\label{sec:Fazit}

\subsection{Soll-/Ist-Vergleich}
\label{sec:SollIstVergleich}

Die im Projektantrag veranschlagten 40~Stunden sind insgesamt
eingehalten worden. Die größten Abweichungen ergaben sich in der
Installations- und der Integrationsphase: Die Grundinstallation
des \ac{PBS} fiel mit ca.\ 5~Stunden etwas aufwändiger aus als
geplant, dafür liefen LDAP-Anbindung und PVE-Integration mit jeweils
ca.\ 30~Minuten deutlich schneller, sodass sich die Differenzen in
Summe ausgeglichen haben. Eine detaillierte Gegenüberstellung der
geplanten und tatsächlichen Stunden je Phase befindet sich in
Tabelle~\ref{tab:Vergleich}.
\tabelle{Soll-/Ist-Vergleich}{tab:Vergleich}{Zeitnachher}

% TODO: Nach Abschluss aller Phasen die endgueltigen Werte in
% Tabellen/Zeitnachher.tex eintragen.


\subsection{Lessons Learned}
\label{sec:LessonsLearned}

Im Laufe des Projekts haben sich vor allem drei Lerneffekte
herausgestellt. Erstens hat sich gezeigt, wie wichtig eine frühe
Abstimmung mit dem Netzwerkteam ist: Der Einbau der \ac{SFP}-
Karte und die \ac{LACP}-Konfiguration waren technisch trivial,
hingen aber von der rechtzeitigen Bereitstellung der Switch-Ports
und der Patchpanel-Verkabelung ab. Zweitens hat der Aufbau auf
Altbestands-Hardware (IBM x3755~M3) klare Grenzen gezeigt, etwa
beim Speicherausbau und bei der Auswahl des Dateisystems
(Hardware-\ac{RAID} statt eines Software-Stacks). Drittens wurde
deutlich, dass eine sauber gepflegte Betriebsdokumentation -- in
diesem Fall im DokuWiki der \ac{ODCF} -- für ein Sysadmin-Projekt
mindestens so wertvoll ist wie die eigentliche Implementierung,
weil der laufende Betrieb anschließend ohne den Prüfling
weiterläuft.

% TODO: ggfs. eigene Erfahrungen ergaenzen (Kommunikation,
% Tooling, Zeitmanagement).


\subsection{Ausblick}
\label{sec:Ausblick}

Mit dem fertigen Blueprint stehen mehrere Folgeschritte offen.
Kurzfristig kann das \ac{ODCF}-Team auf Basis der erstellten
Referenzkonfiguration eine produktive \ac{PBS}-Instanz
aufsetzen und die lokalen Hypervisor-Backups schrittweise auf den
\ac{PBS} migrieren. Mittelfristig sind eine Off-Site-Replikation
des Datastores sowie die langfristige Ablösung der \ac{TSM}-/Tape-
Sicherung als zusätzliche Sicherungsebene denkbar. Darüber hinaus
lässt sich das eingerichtete Monitoring um weitere Metriken und
Alerts ausbauen (z.\,B.\ Trend-Alerts auf Datastore-Belegung
oder Verify-Fehlerquoten), sodass sich der Betrieb mittelfristig
nahezu vollständig automatisiert überwachen lässt.

% TODO: ggfs. konkretere Folgeprojekte mit Frank abstimmen und
% hier referenzieren.
